%% The first command in your LaTeX source must be the \documentclass command.
%%
%% Options:
%% twocolumn : Two column layout. Do not use twocolumn for papers submitted to CEUR-WS!
%% hf: enable header and footer.
\documentclass[
% twocolumn,
% hf,
]{ceurart}

%%
%% One can fix some overfulls
\sloppy

%%
%% Minted listings support 
%% Need pygment <http://pygments.org/> <http://pypi.python.org/pypi/Pygments>
\usepackage{xcolor}
\usepackage{listings}

% --- JavaScript + Unravel-RDF DSL ------------------------------------
\lstdefinelanguage{JavaScript}{
  keywords={const, let, var, function, return, if, else, for, while,
            class, new, await, async},
  sensitive=true,
  comment=[l]{//},
  morecomment=[s]{/*}{*/},
  morestring=[b]",
  morestring=[b]'
}

\lstdefinestyle{unraveljs}{
  language=JavaScript,
  basicstyle=\ttfamily\small,
  % Mots-clés JS (const, etc.) en bleu
  keywordstyle=\color{blue}\bfseries,
  commentstyle=\color{gray}\itshape,
  stringstyle=\color{red!60!black},
  numbers=none,
  frame=single,
  breaklines=true,
  columns=fullflexible,
  % Mots-clés supplémentaires pour mettre en évidence l'API Unravel
  morekeywords={UnravelSession,something,from,out},
  % API Unravel en teal gras
  keywordstyle=[2]\color{teal}\bfseries,
  classoffset=1,
  otherkeywords={UnravelSession,something,from,out},
  classoffset=0,
  % Autoriser l'injection de LaTeX dans le code JS
  escapeinside={(*}{*)}
}

% --- SPARQL -----------------------------------------------------------
\lstdefinelanguage{SPARQL}{
  morekeywords={
    SELECT, CONSTRUCT, ASK, DESCRIBE,
    WHERE, PREFIX, BASE,
    FILTER, OPTIONAL, GRAPH, UNION,
    LIMIT, OFFSET, DISTINCT, REDUCED,
    ORDER, BY, GROUP, HAVING
  },
  sensitive=true,
  comment=[l]{\#},
  morestring=[b]"
}

\lstdefinestyle{sparqlstyle}{
  language=SPARQL,
  basicstyle=\ttfamily\small,
  keywordstyle=\color{blue}\bfseries,
  commentstyle=\color{gray}\itshape,
  stringstyle=\color{red!60!black},
  numbers=none,
  frame=single,
  breaklines=true,
  columns=fullflexible,
  escapeinside={(*}{*)}
}

% to active note:
%\usepackage{todonotes}
% to unactive note:
\usepackage[disable]{todonotes}

%%
%% end of the preamble, start of the body of the document source.
\begin{document}

%%
%% Rights management information.
%% CC-BY is default license.
\copyrightyear{2022}
\copyrightclause{Copyright for this paper by its authors.
  Use permitted under Creative Commons License Attribution 4.0
  International (CC BY 4.0).}

%%
%% This command is for the conference information
\conference{ISWC 2026}

%%
%% The "title" command
\title{Unravel-RDF: A JavaScript Framework for Building Web Applications over RDF Knowledge Graphs}

\tnotemark[1]

\author[1]{Olivier Filangi}[%
  orcid=0000-0002-2094-3271,
  email=olivier.filangi@inrae.fr,
]

\cormark[1]
\address[1]{Institute for Genetics, Environment and Plant Protection (IGEPP), National Research Institute for Agriculture, Food and Environment (INRAE), Université Rennes, Institut Agro, Le Rheu, France}

\author[2]{Nils Paulhe}[%
  orcid=0000-0003-4550-1258
]

\address[4]{Vrije Universiteit Amsterdam, De Boelelaan 1105, 1081 HV Amsterdam, The Netherlands}

\author[3]{Guillaume Laisney}[%
  orcid=0009-0005-5653-3257
]

\address[5]{University of Skövde, Högskolevägen 1, 541 28 Skövde, Sweden}

\author[4]{Clément Frainay}[%
  orcid=0000-0003-4313-2786
]

\author[5]{Franck Giacomoni}[%
  orcid=0000-0001-6063-4214
]

%% Footnotes
\cortext[1]{Corresponding author.}


\begin{abstract}
Semantic Web technologies are widely used to integrate heterogeneous
life science resources as RDF knowledge graphs exposed via SPARQL
endpoints. While these graphs support rich querying, exploiting them
in interactive web applications remains difficult for frontend
developers. Existing tools focus on data access and query execution,
but provide limited support for application-oriented abstractions
such as composable query builders or domain-specific languages
tightly aligned with UI logic.

We present Unravel-RDF, a JavaScript framework that introduces a
focus-based domain-specific language (DSL) for constructing SPARQL
queries together with a composable API tailored to event-driven user
interfaces. Unravel-RDF offers functional abstractions for incremental
query construction, asynchronous execution, pagination, and caching,
while delegating query processing to existing SPARQL engines. We
evaluate Unravel-RDF in the FORVM ecosystem, which interconnects
biomedical and plant knowledge graphs through shared chemical
entities, and demonstrate a web application for cross-domain
exploration of metabolomics relationships across multiple SPARQL
endpoints. This case study suggests that Unravel-RDF can bring
knowledge graph querying closer to modern frontend development
practices and improve the modularity and responsiveness of
metabolomics exploration tools.
\end{abstract}


\begin{keywords}
  Knowledge graphs \sep
  JavaScript framework \sep
  Frontend development \sep
  Domain-specific language (DSL) \sep
  Metabolomics \sep
\end{keywords}

\maketitle

\section{Introduction}

Semantic Web technologies are widely used to represent and query knowledge in the life sciences\citep{wu_semantic_2014}. 
Many resources are published as RDF knowledge graphs and made available to the community via downloadable datasets 
or SPARQL endpoints, thereby enabling interoperability between heterogeneous datasets describing molecular entities, 
biological processes, pathways, and phenotypic characteristics across diverse organisms\citep{wilkinson_fair_2016,LODCloudLifeSciences}.
In bioinformatics, these approaches support the integration and reuse of theses resources, but their effective use within analysis platforms 
and web applications still poses significant technical challenges.

Developers building web frontends on top of RDF knowledge graphs face a fundamental challenge: 
bridging low-level RDF graph models with application logic. Frontend developers typically 
either embed raw SPARQL queries directly in their code or rely on low-level RDF/JS-based libraries that expose 
graph primitives and query execution but offer limited support for structuring application data access.

The RDF JavaScript Libraries Community Group - Tools\footnote{\url{https://www.w3.org/groups/cg/rdfjs/tools/}}, including engines such as Comunica\citep{taelman_comunica_2018}, has been highly successful at standardizing RDF representations and query interfaces 
across heterogeneous data sources. However, these tools are intentionally designed as foundational building blocks, and do not provide the kind 
of higher-level, application-oriented abstractions—such as composable query builders, domain-specific languages (DSLs), or data mappers—that 
frontend developers routinely rely on in relational database ecosystems.

As a result, there remains a significant gap between RDF query formulation and modern frontend development practices, 
leading to brittle integrations and weak alignment between knowledge graph queries and application-level structures.

To address this gap, we introduce Unravel-RDF, a JavaScript framework that combines a DSL for SPARQL query construction with a 
composable API tailored to event-driven user interface applications. 
Our design provides functional abstractions for incremental, focus-based query building, asynchronous execution, pagination, and caching 
over RDF data sources. 
We evaluate Unravel-RDF on the FORVM DiseaseChem knowledge graph, which integrates associations between chemical compounds and biomedical 
concepts, as well as on the experimental FORVM Plants graph, built from analogous associations applied to plant-domain ontologies derived 
from the Planteome project. To this end, we implement an inter-domain exploration workflow, deployed as a web application, which leverages 
these abstractions to enable interactive navigation between concepts across multiple SPARQL endpoints.

\section{Architecture and Programming Model}

\begin{figure}[t]
    \centering
    \includegraphics[width=1.0\linewidth]{Unravel-rdf-architecture.png}
    \caption{Unravel-RDF architecture: a layered design connecting RDF data sources, the Unravel query abstraction layer, and UI integration components.}
    \label{fig:architecture}
\end{figure}

\subsection{Layered Architecture}

Figure~\ref{fig:architecture} illustrates the architecture of
Unravel-RDF as a three-layer system: an \emph{Application Development
Layer}, an \emph{Unravel Query Abstraction Layer}, and an \emph{RDF Data
\& Query Layer}. The Application Development Layer hosts
JavaScript/TypeScript code, HTML-based user interfaces, and Node.js
services, together with Unravel sessions that are embedded into web
applications. This layer includes hooks and UI components (e.g., built
with D3.js, React, or Vue) that consume Unravel queries, as well as
three immutable objects, \texttt{UnravelConfiguration},
\texttt{UnravelSession}, and \texttt{UnravelQuery}, which encapsulate
configuration, query construction, and execution respectively.

The Unravel Query Abstraction Layer is implemented in Scala/Scala.js and
provides the core query model. It is organized into several functional
subsystems: a configuration module, a core query DSL, a high-level query
API, UI integration utilities, and a variable scoping mechanism.
Together, these components offer a composable, focus-based abstraction
over SPARQL queries that can be safely reused across application code.
A dedicated Comunica facade\footnote{https://www.scala-js.org/libraries/facades.html} 
connects this abstraction layer to the JavaScript RDFJS ecosystem and exposes a typed interface to the
underlying SPARQL engines.

Finally, the RDF Data \& Query Layer corresponds to the existing
JavaScript/TypeScript ecosystem for RDF and SPARQL. In our deployment,
this layer is instantiated by a Comunica SPARQL engine that can query
over multiple RDF sources, including HTTP-accessible SPARQL endpoints
and local or remote RDF files. Unravel-RDF therefore focuses on bridging
application-level query construction and UI logic with these
heterogeneous RDF data sources, without reimplementing query processing.

\subsection{Programming Model}

At its core, Unravel-RDF provides a domain-specific language (DSL) for
constructing SPARQL queries. Instead of relying on raw query strings,
the DSL enables the composition of structured query fragments.
Query construction is incremental and follows a lambda-closure style,
where functions capture context and enable modular composition.
All manipulated objects are immutable, enforcing a functional
programming paradigm that improves the composability of query
transformations.

Queries are constructed relative to a focus variable, which serves as
an implicit context for navigating RDF graphs. Focus-binding primitives
(\texttt{something}, \texttt{from}) introduce or rebind this variable,
defining the scope of query construction. In contrast, graph-expansion
primitives (\texttt{out}, \texttt{in}, \texttt{traverse}) add triple
patterns relative to the current focus while updating it to the newly
reached node. These primitives accept both RDF IRIs and SPARQL variables
as arguments, so the same combinators can express schema-level templates
and data-level queries in a uniform way. Table~\ref{tab:query-primitives}
summarizes the core primitives and their semantics, and the following
example shows how they can be composed to build a query centered on a
compound entity acting as the focus variable. In this example\ref{fig:unravel-sparql-example}, 
\texttt{session2} extends an initial query over the focus variable \texttt{compound}, 
illustrating how Unravel-RDF supports incremental refinement of SPARQL queries.
In practice, these primitives are exposed through \texttt{UnravelSession} objects, which provide a fluent 
API for constructing and reusing focus-based queries.

\begin{table}[t]
\centering
\setlength{\tabcolsep}{3pt}
{\footnotesize
\caption{Core query building primitives in Unravel-RDF}
\label{tab:query-primitives}
\begin{tabular}{|l|p{3cm}|p{6cm}|}
\hline
\textbf{Primitive} & \textbf{Role} & \textbf{Description} \\
\hline
\texttt{something(name,f)} 
  & Focus introduction 
  & Introduces a fresh SPARQL variable as focus; \texttt{f} builds patterns relative to it. \\
\hline
\texttt{from(name,f)} 
  & Focus rebinding 
  & Reuses an existing variable as focus; \texttt{f} continues query construction from it. \\
\hline
\texttt{out(prop,obj,f)} 
  & Forward expansion 
  & Adds \texttt{<focus> prop obj} and sets focus to \texttt{obj}. \\
\hline
\texttt{in(prop,subj,f)} 
  & Backward expansion 
  & Adds \texttt{subj prop <focus>} and sets focus to \texttt{subj}. \\
\hline
\texttt{traverse(prop,node,f)} 
  & Bidirectional 
  & Adds \texttt{\{<focus> prop node\} UNION \{node prop <focus>\}} and sets focus to \texttt{node}. \\
\hline
\end{tabular}
}
\end{table}

\begin{figure}[t]
\centering
\caption{Incremental query construction with Unravel-RDF and resulting SPARQL}
\label{fig:unravel-sparql-example}

{\footnotesize
\begin{minipage}[t]{0.50\textwidth}
\centering
\begin{lstlisting}[style=unraveljs,basicstyle=\footnotesize\ttfamily]
const session =
  UnravelSession(config)
    .something("(*\textcolor{purple}{compound}*)",
      (*\textcolor{purple}{compound}*) =>
        (*\textcolor{purple}{compound}*).out("(*\textcolor{purple}{?rel}*)","(*\textcolor{green!50!black}{CHEBI:24279}*)"));

const session2 =
  session.from("(*\textcolor{purple}{compound}*)",
    (*\textcolor{purple}{compound}*) =>
      (*\textcolor{purple}{compound}*).out("(*\textcolor{green!50!black}{rdf:type}*)","(*\textcolor{purple}{?type}*)"));
\end{lstlisting}
\end{minipage}%
\hspace{0.02\textwidth}%
\begin{minipage}[t]{0.44  \textwidth}
\centering
\begin{lstlisting}[style=sparqlstyle,basicstyle=\footnotesize\ttfamily]
SELECT (*\textcolor{purple}{?compound}*) (*\textcolor{purple}{?rel}*) (*\textcolor{purple}{?type}*)
WHERE {
  (*\textcolor{purple}{?compound}*) (*\textcolor{purple}{?rel}*) (*\textcolor{green!50!black}{CHEBI:24279}*) .
  (*\textcolor{purple}{?compound}*) (*\textcolor{green!50!black}{rdf:type}*) (*\textcolor{purple}{?type}*) .
}
\end{lstlisting}
\end{minipage}
}
\end{figure}

\subsection{Execution Backend}

Unravel-RDF delegates SPARQL query execution to existing engines via 
a Scala.js Comunica facade. This facade maps UnravelQuery objects 
onto RDFJS-compatible data sources and SPARQL endpoints, exposing 
a typed interface to the JavaScript query engine.

Through Comunica’s federation capabilities, Unravel-RDF can query heterogeneous, 
distributed RDF datasets, while leaving query planning, optimization, and low-level 
data access entirely to the underlying engines.

\subsection{User Interface Integration}

On the application side, Unravel-RDF is designed to integrate with
modern, event-driven user interfaces. The execution model is based on
JavaScript promises combined with a subscription mechanism, enabling
asynchronous and staged processing of queries. Results are delivered in
batches through a pagination mechanism and exposed as lazy pages
(windows of results) that are only fetched when requested, which avoids
long-running SPARQL requests and reduces perceived latency.

These abstractions are intended to be consumed directly from
hook-based and component-based UI frameworks such as React, Vue, or
D3-powered visualizations. Application code can subscribe to
\texttt{UnravelQuery} objects, trigger incremental loading of result
pages, and update views as data becomes available, following typical
reactive programming patterns. Additionally, a caching layer prevents
redundant query execution by reusing previously retrieved results,
thereby improving performance and responsiveness. Overall, this
integration brings SPARQL-based querying closer to modern frontend
development practices while preserving a clear separation between UI
concerns, query abstraction, and engine-level execution.

\subsection{Implementation and Distribution}

Unravel-RDF is implemented in functional Scala and compiled to
JavaScript using Scala.js, which allows the same core query abstraction
layer to be reused across browser-based and Node.js applications. It is
distributed through standard web channels, including CDN delivery and
the npm ecosystem, facilitating integration into modern JavaScript and
TypeScript build pipelines. The source code and documentation are
publicly available in the project repository\footnote{\url{https://forge.inrae.fr/p2m2/unravel-rdf}},
enabling reuse and extension by third-party developers.

\section{Case Study: FORVM}

We evaluate Unravel-RDF in the context of the FORVM ecosystem, which
provides a cross-domain knowledge graph infrastructure spanning
biomedical and plant sciences. The biomedical component, \textit{FORVM
ChemDisease}, is a curated knowledge graph derived from statistical
co-occurrence of MeSH-annotated chemical and disease concepts in PubMed
literature, while \textit{FORVM Plants} extends this approach to plant
sciences by leveraging Planteome ontologies to represent plant
developmental stages, diseases, and associated traits as ontological
concepts.

These two knowledge graphs are interconnected through shared chemical
compounds, enabling cross-domain exploration of metabolomics
relationships across multiple SPARQL endpoints, including FORVM ChemDisease
and FORVM Plants, as well as RDF Portal SPARQL endpoints\citep{kawashima_nbdc_2018}
for resources such as ChEBI\citep{de_matos_chebi_2010} and PubChemRDF\citep{fu_pubchemrdf_2015}, 
and MeSH-annotated concepts from PubMed\citep{coletti_medical_2001}.

We implement an inter-domain exploration workflow as a web application
built with Unravel-RDF, in which the interface displays clusters of
nodes linking chemical compounds to both MeSH disease terms and
Planteome concepts, revealing transitive connections between biomedical
and plant knowledge.

In this application, Unravel-RDF orchestrates compositional query
construction and reuse of focus-based query fragments while issuing
separate SPARQL queries to the different endpoints as needed. The
resulting query logic is encapsulated in reusable Unravel-RDF sessions
and queries, so that the UI can trigger navigation steps (e.g.,
expanding a compound to related diseases or traits, then following links
to external knowledge graphs) without embedding raw SPARQL strings in
the frontend code. The case study UI is available at
\url{http://kg.p2m2.inrae.fr/}.

\begin{acknowledgments}
We gratefully acknowledge support from the French national infrastructure 
MetaboHUB, funded by the French National Research Agency (ANR) under the 
France~2030 programme, for the development of the Metabolomics Semantic 
Data Lake e-infrastructure\cite{MetabolomicsSemanticDataLake}. This platform 
enables the construction of the FORVM knowledge graphs and the application 
layers for their exploitation, from which Unravel-RDF directly benefits.

We also thank INRAE’s Digit-Bio metaprogramme for funding the SEED (Semantic 
Engineering for Enhanced Data in Plant Metabolomics) project, which 
extends the FORVM methodology to plant metabolomics and food processing 
domains, and supports the further development of the Unravel-RDF DSL 
for leveraging this new body of knowledge.
\end{acknowledgments}

\section*{Declaration on Generative AI}

During the preparation of this work, the authors used Perplexity and
ChatGPT in order to: paraphrase and reword text, improve writing
style and clarity, and perform grammar and spelling checks. After using these
tools, the authors reviewed and edited the content as needed and take
full responsibility for the publication’s content.


%%
%% Define the bibliography file to be used
\bibliography{unravel-biblio,static-extra}

\appendix

%\section{Supplementary Material}

%\begin{figure}[ht]
%\centering
%\includegraphics[width=\linewidth]{forvm-unravel-screenshot.png}%
%\caption{Screenshot of the FORVM use case web interface built with Unravel-RDF, showing clusters of compounds, MeSH disease terms, and Planteome concepts.}
%\label{fig:forvm-screenshot}
%\end{figure}

\section{Online Resources}

The following resources provide additional material related to the
FORVM ecosystem and the Unravel-RDF programming model:
\begin{itemize}
  \item Unravel-RDF source code and documentation:
    \href{https://forge.inrae.fr/p2m2/unravel-rdf}{https://forge.inrae.fr/p2m2/unravel-rdf}
  \item FORVM DiseasesChem :
    \href{https://forum-webapp.semantic-metabolomics.fr}{https://forum-webapp.semantic-metabolomics.fr}
  \item Supplementary examples and usage tutorials:
    \href{https://example.org/unravel-examples}{https://example.org/unravel-examples}
\end{itemize}

\end{document}
%%
%% End of file
